\chapter{Reading the benchmarks honestly}
\label{ch:benchmarks}

\newthought{Long-context benchmarks} have a credibility problem in 2026.
Vendors report numbers that look spectacular; independent reproductions
either don't exist or contradict the vendor numbers; the underlying
evaluation setups (which model, which prompt template, which scoring
method, which context window) often differ from one paper to the next.

This chapter is a survival guide for reading SubQ's published numbers
without either dismissing them or being sold by them.\sidenote{Source for
all of SubQ's numbers in this chapter: \citet{subq2026ssa}. The
``third-party verified'' qualification on these numbers refers to a
specific evaluation lab; the full methodology was not public as of the
last update of this book.}

\section{The three benchmarks that matter}

\paragraph{RULER.} \citet{hsieh2024ruler} is the de-facto long-context
benchmark in 2026. A battery of 13 sub-tasks covering needle-in-a-haystack,
multi-key retrieval, multi-hop reasoning, and aggregation, all at
configurable context lengths from 4K to 128K. Reported as average
accuracy across the 13 tasks at a given length.

\paragraph{MRCR v2.} \citet{chen2025mrcr} is the harder successor. Multi-
document reasoning + retrieval. Tests whether a model can read multiple
related documents at long context and answer a question requiring
cross-document inference. The 2026 version includes a leaderboard.

\paragraph{Wall-clock prefill speedup.} The non-quality benchmark: how
fast does the model encode a long prompt? Usually reported relative to a
baseline (FlashAttention-2 on a B200, in SubQ's case).

\section{What SubQ reports}

From \citet{subq2026ssa}, paraphrased:

\begin{center}
\begin{tabular}{ll}
\toprule
Metric & SSA / SubQ result \\
\midrule
RULER \@ 128K (avg of 13 tasks) & \(\sim\) 95\% \\
Prefill speedup vs FA-2 \@ 1M tokens (B200) & \(\sim\) 52$\times$ \\
Attention-FLOP reduction vs dense \@ 1M & \(\sim\) two orders of magnitude \\
MRCR v2 (overall) & 65.9\% \\
\bottomrule
\end{tabular}
\end{center}

\paragraph{What the RULER number means.} 95\% on RULER \@ 128K is
genuinely strong. Frontier proprietary models (Claude 3.5/4, GPT-4/5,
Gemini) report numbers in the 90--96\% range at this length; SubQ being
competitive here is the headline.

\paragraph{What the prefill speedup means.} 52$\times$ vs FlashAttention-2
on a B200, at 1M tokens. The baseline matters: FA-2 is itself an
aggressive optimisation of dense attention. The speedup is for
\emph{prefill} specifically — the initial encoding of the prompt — which
is where the \(n^2\) hits hardest. Decode (per-token generation) speedup
will be smaller because both architectures hit the KV-cache memory wall
in similar ways.

\paragraph{What the MRCR v2 number means — the part the marketing
underplays.} The honest read of the MRCR v2 leaderboard at the time of
writing (paraphrased from \citealp{subq2026ssa}'s own table):

\begin{center}
\begin{tabular}{lr}
\toprule
Model & MRCR v2 \\
\midrule
Claude Opus 4.6     & 78.3\% \\
GPT-5.5             & 74.0\% \\
\textbf{SubQ (SSA)} & \textbf{65.9\%} \\
GPT-5.4             & 36.6\% \\
Claude Opus 4.7     & 32.2\% \\
Gemini 3.1 Pro      & 26.3\% \\
\bottomrule
\end{tabular}
\end{center}

SubQ is \emph{in the conversation} but \emph{not at the top}. The frontier
proprietary models have a clear lead. The relative position is what
matters: at a presumably much smaller training-compute footprint, SSA
is competitive with — but not yet beating — the best dense transformers
at multi-document long-context reasoning. That is a credible architectural
result.

\keyidea{Long-context benchmarks are a moving target. SubQ's strongest
honest claim is that SSA is competitive with frontier proprietary
transformers on RULER and MRCR v2 \emph{while running 50$\times$ faster
on prefill at 1M tokens}. ``Faster than the best, comparable in quality''
is a much stronger claim than ``best at everything''.}

\section{What is missing}

A short list of things the public materials do not (yet) settle:

\begin{itemize}
\item Training data composition, total training tokens, and parameter
count. Without these, comparing to GPT-5 or Claude Opus 4.6 isn't
apples-to-apples.
\item The full text of the third-party verification.
\item Decoding wall-clock numbers (the prefill speedup is for
\emph{prompt encoding}; per-token decode performance is a separate
question).
\item Behaviour at $> 1$M tokens — RULER \@ 1M, MRCR v2 \@ 1M.
\end{itemize}

The right posture is ``cautiously impressed and waiting for
reproducibility''. That posture is the one this book has tried to model
throughout.
